\chapter{Theoretical Background}\label{chapter:theoretical_background}

  \section{Medical foundations of EEG signals}

  \section{Brain rythms}

  \section{Recording and measurement of EEG signals}
    % Mention common placements
  \section{Denoising and preprocessing of EEG signals}
    After acquiring the EEG signals, one has to follow further steps before performing a classifaction task on it. Since \ac{EEG} is a non-invasive technique, other signals, so called artifacts, are measured as well. Before discussing possible techniques to remove those data, different types of artifacts are presented with examples included. 
    \subsection{Types of artifacts in EEG signals}
        \begin{table}[h]\label{tab:artifacts}
        \begin{tabular}{|p{0.45\textwidth}|p{0.45\textwidth}|}
          \hline
          \textbf{Physiological artifacts} & \textbf{Non-Physiological artifacts} \\
          \hline
          Ocular-related artifacts: eye movements and blinking & Power line interference \\
          \hline
          \ac{EMG} artifacts in head movements & Poor placement of electrode on scalp \\
          \hline
          \ac{ECG} artifacts in heart beat & Malfunction of a part in the EEG recording system \\
          \hline
          Scalp perspiration and movements & Digital artifacts due to loose wiring/circuit board connections\\
          \hline
        \end{tabular}
        \centering
        \caption{Types of artifacts in EEG \cite{hu2019eeg}}
      \end{table}
      Before illustrating what types of artifacts exist, it is important to mention that artifacts themselves do not destroy \ac{EEG} signals. They just add up on top of the neural activity. Therefore, before removing any of this data to the types and characteristics have to be explained thoroughly. In \autoref{tab:artifacts}, an overview of forms of artifacts can be found. In general, artifacts can be put into two different groups, the \textit{Non-Physiological artifacts} and the \textit{Physiological artifacts}. As the naming already suggests, while Non-physiological artifacts are distortions coming from the outside world, Physiological artifacts are intereferences by the human body. \cite{hu2019eeg} Both of these groups will be further described in the following two sections.
      \subsubsection{Physiological artifacts}
        As already metioned, Physiological artifacts have the characteristic of coming from the human body, including ocular movements, head movements, heartbeats and muscle noises. For most of theses artifacts, it is benefitial to record their corresponding Biosignals. In that way, it is possible to retrace those artifacts in the measured signal. \cite{Uriguen_2015}\par
        Ocular artifacts can be measured by \ac{EOG} and are mainly produced by eye movements and blinking. The latter one generally prodces a much larger amplitude than than the \ac{EEG} activity. Because of the anatomy of a human, \ac{EOG} contaminations mostly occur in the frontal electrodes. \cite{Uriguen_2015} \par
        An even more common artifact comes in muscular activity which can be measured in an \ac{EMG} signal. Appearing in studies where the person is awake, the main actors for those kind of contaminations are swallowing, talking and walking, depending on the experiment environment. It is known that mostly cranial \ac{EMG} signals are in charge for disadvantegous effects in the recorded signal. However, it is difficult to detect and correct them as \ac{EMG} not only represents a wide range on the spectral distribution, but also does not repeat making it even more difficult to characterise. \cite{Uriguen_2015} \par
        Last but not least, the activity of the heart (\ac{ECG}) plays a significant role in contamination of measured \ac{EEG} signals. The amplitude of them usually is low but also depends on electrode placing. In contrast to \ac{EMG} data, a regular pattern can be visible, which can be further classified by comparing it to \ac{ECG} measurements. \cite{Uriguen_2015}\par
      \subsubsection{Non-Physiological artifacts}
        On the other hand, there are the so called non-physiological artifacts coming from outside-world intereference. The most common source is hereby the electrical intereference originating from the power line (50 Hz in Europe). This generally can be reduced by shielding cable or moving the subject far away from the power source. \cite{hu2019eeg} \par
        Another very common non-physiological artifact is the poor placement of electrodes on the scalp as the created DC potential is disturbed. Therefore, electrodes should be checked regularly to ensure high quality recordings. \cite{hu2019eeg} \par
    \subsection{Artifact Removal techniques}
      Before presenting state of the art techniques for removing artifacts, it is important to mention that most of these algorithms and techniques concern physiological artifacts. Most of the non-physiological artifacts can be corrected or even completely avoided by controlling the recording environment. By checking the placement of the patient and the electrodes on the scalp, non-physiological remain a small problem in artifact removal compared to handlig physiological artifacts. \cite{hu2019eeg,Uriguen_2015} \par
      The first technique is the so called (Frequency) filtering. It not only removes the power line interference but also removes high- and low-frequency noises. Depending on the goal of the filtering, different filtering techniques are available, mainly low-pass filter, high-pass filter, band-pass filter and band-stop filter. While a low-pass filter removes all frequencies above a chosen threshold, a high-pass filter eliminates frequencies below a boundary. The implementation of band-pass filters on the other requires the definition of a lower and and upper frequency boundary. When applying, it removes all frequencies that are below the lower frequency and above the higher frequency (so keeping all frequencies that are within the given constraints). Last but not least, a band-stop filter does the complete oppisite. It removes all frequencies that are withing the chose interval, thus keeping the frequencies below the lower bound and above the upper bound. Frequency filtering should be the first step of preprocessing \ac{EEG} signals (even before segmentation) in order to remove low- and high-frequency noise outside of the spectrum of brain rythms. \cite{hu2019eeg} \par
      Whent talking about algorithms and techniques dedicated to remove artifacts, assumptions have to be met about the recorded signal. The most important assumption hereby is that the recorded signal is the sum of the brain activity and noise. A more generalized version describes the measured \ac{EEG} signals as linear mixtures of waveforms descending from multiple brain sources and artifacts, ultimately propagating to the scalp. Mathematically, this can be described by following mathematical equation: 
      \begin{equation}
        \mathbf{X} = \mathbf{A}\mathbf{S} + \mathbf{V},
        \label{eq:mixing}
      \end{equation}
      where $\mathbf{X} \in \mathbb{R}^{n \times N}$, $\mathbf{A} \in \mathbb{R}^{n \times m}$, $\mathbf{S} \in \mathbb{R}^{m \times N}$,$\mathbf{V} \in \mathbb{R}^{n \times N}$. While $N$ represents the total count of samples measured, $n$ stands for the amount of \ac{EEG} channels and $m$ for the amount of total sources the signal originates from. In \autoref{eq:mixing}, $\mathbf{X}$ describes the \ac{EEG} signals resulting from the measurement. While $\mathbf{V}$ stands for added noise \footnote{In the following, it is assumed that $\mathbf{V} = 0$}, $\mathbf{A}$ is an unknown mixing matrix and $\mathbf{S}$ is a matrix of unknown Sources representing brain sources and artifacts. \cite{Uriguen_2015} \par
      With this model, the goal is now to solve the \ac{BSS} problem, so in this case seperating neural activity from activity. For this problem, component based methods are used because they find principal or independent components that are equal to the input \ac{EEG} channels. After processing in the transformed domain, an inverse transformation is applied with the corrected components resulting in processed \ac{EEG} data. \cite{Uriguen_2015} \par
      In general, \ac{BSS} aims to estimate $\mathbf{S}$ on the basis of $\mathbf{X}$ in \autoref{eq:mixing}. It has to be mentioned that the effectiveness of those methods lies in different assumptions that have to be met in a certain extent, e.g. independence or non-Gaussanity. The more fulfilled these assumptions are, the more effective the \ac{BSS} method is. Apart from that, the success of such methods also depends on the correct identification of the sources as either brain activity or artifacts, which can serve as a challenge depending on the extent of artifacts. \cite{Uriguen_2015} \par
      One of those methods is \ac{PCA}, which makes use of an orthogonal transformation to turn the observed data ot linearly uncorrelated variables (Principal components). The resulting components have the property of having the largest possible variances while being orthogonal to each other. However, current research suggests that artifacts and cerebral activity can be modeled better as independent instead of orthogonal. To make use of this method \ac{ICA} has emerged as the leading method for \ac{BSS}. \cite{Uriguen_2015} \par
      \ac{ICA} consists of the use of multiple related methods seperating linearly mixed signals by taking advantage of the statistical independence of the sources. In most cases, the algorithm work based on \ac{HOS}. After prewhitening the covariance matrix of the data, \ac{ICA} tries to find a linear transformation $\mathbf{W}$ for estimated sources $\hat{\mathbf{S}} = \mathbf{W} \mathbf{X}$. $\hat{\mathbf{S}}$ is then supposed to be as independent as possible. How effective the \ac{ICA} not only depends on the independence of sources but also the full column rank of the mixing matrix $\mathbf{A}$. Since this algorithm relies on statistical features, it is important that the amount of samples is sufficiently large. In most cases, it is recommended to use all available data under the condition that cerebral activity alongside with the artifacts are spatially stationary through time. \cite{Uriguen_2015} \par

  \section{Feature extraction and signal decomposition of EEG signals}